\section{Periodicity and the uniform bound}\label{sec:consequences}
The finite sieve theorem will be proved independently of the component
bound. The following conditional reduction explains why that theorem
suffices, and records the contribution of the finitely many Gaussian
primes removed by the sieve. The periodicity principle also appears in
Vardi~\cite[Section 6, Proposition 6.2]{Vardi}.

\begin{proposition}[From a periodic obstruction to a uniform bound]
\label{prop:periodic}
Fix $D\ge1$ and a finite set $\mathcal P$ of rational primes congruent
to $1$ modulo $4$. Suppose that $\mathcal A(\mathcal P)$ contains no
infinite self-avoiding walk with steps at most $D$. Put
\[
 Q=\prod_{p\in\mathcal P}p,\qquad
 K_D=\#\{v\in\Z[i]:|v|\le D\},
\]
and let $E$ be the set of associates of the selected Gaussian factors.
For every component $C$ of $G_D$,
\[
 \#C\le\max\{Q^2,\ |E|+(K_D-1)|E|Q^2\}.
\]
\end{proposition}
\begin{proof}
The graph on $\mathcal A(\mathcal P)$ with this step bound is locally
finite and invariant under translations by $Q\Z^2$. An infinite
component would contain an infinite simple path: the tree of finite
simple paths from one vertex is finitely branching, has vertices at
arbitrarily large depths, and one can successively choose a child with
arbitrarily deep descendants. Hence every component is finite.

Suppose that two distinct points $x,y$ of a component differ by
$v\in Q\Z^2$. Translation by $v$ sends that component to a component
containing $y$, so sends it to itself. A finite nonempty subset of
$\Z^2$ cannot be invariant under a nonzero translation. Each component
therefore injects into $(\Z/Q\Z)^2$ and has size at most $Q^2$.

A Gaussian prime divisible by a selected factor is an associate of that
factor. Thus, after deleting $E$, the Gaussian-prime graph is a subgraph
of the avoiding graph. A component meeting $E$ can meet at most
$(K_D-1)|E|$ components of the graph with $E$ removed: each such
component has an edge to $E$, whose total degree is at most that number.
Adding back $E$ gives the claimed bound.
\end{proof}


Applying Proposition~\ref{prop:periodic} to the set supplied by
Theorem~\ref{thm:sieve} proves Theorem~\ref{thm:main} for $D\ge1$;
the case $D<1$ was settled in the introduction. For the weaker
infinite-walk conclusion, one can simply delete an initial segment
containing the finitely many exceptional Gaussian primes: the remaining
tail would lie in $\mathcal A(\mathcal P_D)$.

Periodicity is what makes the bound uniform in the starting point.
Absence of an infinite path in a general locally finite graph does not
bound the sizes of its finite components. We now turn to the arithmetic,
entropy, and geometry needed to produce the finite obstruction.
